\documentclass[12pt]{letter}
\usepackage[utf8]{inputenc}
\usepackage[T1]{fontenc}
\usepackage{libertine}
\usepackage[margin=1in]{geometry}
\usepackage{hyperref}

\signature{Jonathan Mace\\Microsoft Research, Redmond, WA\\[4pt]
Brian R.\ Mace\\University of Auckland, New Zealand}

\address{Jonathan Mace\\Microsoft Research\\Redmond, WA 98052\\United States\\
jonathanmace@microsoft.com}

\date{\today}

\begin{document}

\begin{letter}{The Editor\\
Proceedings of the Royal Society A:\\
Mathematical, Physical and Engineering Sciences\\
The Royal Society\\
6--9 Carlton House Terrace\\
London SW1Y 5AG\\
United Kingdom}

\opening{Dear Editor,}

We wish to submit the enclosed manuscript, \emph{Geometry Mediates Excitation,
Resonance, and Identifiability in Fluid-Filled Viscoelastic Shells}, for
consideration as a Research Article in \emph{Proceedings of the Royal Society~A}.

The paper advances a unifying thesis for the mechanics of fluid-filled
viscoelastic shells: shell geometry simultaneously filters the forcing field,
organises the modal spectrum, and determines the invertibility of the
corresponding parameter-estimation problem.  These three roles --- excitation
filtering, spectral organisation, and inverse identifiability --- have
previously been studied in isolation across disparate application domains.  The
present work establishes them as logically connected consequences of the same
geometric structure.

The manuscript contributes five formal results within a reduced analytical
framework:
\begin{enumerate}
  \item \textbf{Theorem (Rank Collapse).} Equivalent-radius reductions of
    non-spherical fluid-filled shells are generically rank deficient: collapsing
    independent geometric parameters into a single scalar size descriptor
    confounds geometry and material stiffness, rendering independent parameter
    recovery impossible.

  \item \textbf{Proposition (Identifiability Lifting).} Oblate asphericity lifts
    this rank deficiency by introducing curvature-dependent mode splitting that
    provides the additional spectral information needed for locally identifiable
    inversion.

  \item \textbf{Proposition (Near-Spherical Asymptotics).} The transition to the
    spherical limit follows a regular expansion,
    $\sigma_3(\varepsilon) = \lambda_1 \varepsilon^2 + O(\varepsilon^4)$,
    confirming that the exact sphere is a singular point in the inverse problem
    and that any previously reported finite identifiability floor is a
    discretisation artefact.

  \item \textbf{Proposition (Forward $\neq$ Inverse Adequacy).} A counterexample
    demonstrates that forward-model adequacy does not imply inverse-model
    adequacy: an equivalent-sphere reduction may predict resonant frequencies
    within engineering tolerance yet fail catastrophically when used for
    parameter estimation.  The structural mechanism is a sensitivity gap along
    the kernel of the geometric reduction.

  \item \textbf{Proposition (Breit--Wigner Absorption).} The energy fraction
    absorbed by the shell from an incident acoustic field follows a
    Breit--Wigner resonance profile, providing a rigorous basis for the
    excitation-filtering role of geometry.
\end{enumerate}

The formal results are supported by cross-application evidence drawn from three
distinct physical systems: the human abdomen (infrasound resonance and
whole-body vibration coupling), the watermelon (acoustic ripeness assessment via
shell stiffness gradients), and the urinary bladder (fill-dependent resonance
tracking).  Additional supporting evidence spans gas-pocket transduction,
scaling-law analysis, borborygmi modelling, and sub-bass perception pathways.

This capstone paper synthesises the results of nine companion papers developed
within the Browntone research programme, each of which addresses a specific
aspect of the geometry--spectrum--inference nexus.  We believe the unifying
framework, the formal inverse-problem results, and the breadth of
cross-application evidence make the manuscript well suited to the
interdisciplinary readership of \emph{Proceedings of the Royal Society~A}.

The manuscript has not been published previously and is not under consideration
elsewhere.  All authors have approved the manuscript and agree with its
submission to \emph{Proceedings of the Royal Society~A}.

We confirm that all necessary ethical approvals have been obtained.  The study
is purely analytical and computational and involves no human or animal subjects.

We have included suggested reviewers with expertise spanning inverse problems
in structural dynamics, shell vibration, spectral geometry, and biomedical
acoustics.

\closing{Yours sincerely,}

\end{letter}
\end{document}
